\chapter*{Kết luận}
\addcontentsline{toc}{chapter}{Kết luận}

Qua quá trình khảo sát bài toán thực tiễn, lựa chọn công nghệ, thiết kế kiến trúc, hiện thực mã nguồn và kiểm thử đánh giá, nhóm chúng em đã hoàn thiện nền tảng IoT nhà kính thông minh GreenHub theo đúng mục tiêu ban đầu. Thay vì xây dựng một mạch điều khiển nhúng rời rạc với ngưỡng cố định, hệ thống được định vị và tổ chức như một nền tảng IoT chuyên biệt hoàn chỉnh, tự chủ 100\% mã nguồn với 27.049 dòng code xuyên suốt từ tầng cảm biến -- chấp hành, tầng truyền thông bảo mật, tầng nền tảng dữ liệu và bộ luật điều khiển cho tới tầng học máy hỗ trợ quyết định canh tác và giao diện Web Console PWA.

Về mặt kỹ thuật, những kết quả nổi bật mà nhóm đạt được thể hiện ở ba phương diện chính. Thứ nhất là cơ chế đảm bảo an toàn đa tầng: từ bộ lọc loại gai nhiễu, lọc trung vị 3 mẫu kết hợp EMA tại Gateway, giới hạn thời lượng chạy, tự ngắt đúng hạn cho 3 kênh relay và failsafe cục bộ khi mất mạng, cho đến kiểm tra chất lượng dữ liệu 5 nhãn, vòng điều khiển hysteresis dựa trên các chỉ số nông học VPD/DLI kèm bảo vệ chống tác dụng phụ VPD và kiểm soát chế độ \texttt{auto}/\texttt{manual} có thời hạn trên máy chủ. Thứ hai là khả năng thống nhất trọn vẹn giữa mô hình nhà kính mini vật lý và công cụ giả lập đa trang trại qua chung một luồng cấp phát HTTP provisioning, MQTT TLS 1.2+ kèm Dynamic Security ACL, cô lập đa tenant theo \texttt{farm\_id} với mô hình quyền \texttt{owner} thống nhất và hệ thống cảnh báo 2 lớp (In-app WebSocket kết hợp còi Buzzer tại chỗ). Thứ ba là phân hệ học máy thực dụng và minh bạch: một mô hình LightGBM toàn cục học từ dữ liệu khí hậu liên nhà kính phục vụ đồng thời hai bài toán dự báo năng suất và gợi ý top-3 cây trồng kèm lời giải thích tiếng Việt, hoàn toàn tách rời khỏi vòng điều khiển thiết bị. Kết quả đo đạc cho thấy toàn bộ 58/58 ca kiểm thử đơn vị Python vượt qua 100\%, mã nguồn giao diện TypeScript không còn lỗi kiểm tra kiểu tĩnh, và mô hình LightGBM phiên bản \texttt{v2} (gộp dữ liệu tổng hợp cùng bộ dữ liệu thực nghiệm WUR) đạt \(R^2 = 0,9387\), \(\text{MAE} = 0,6792\;\text{kg/m}^2\) cùng \(\text{NDCG@3} = 0,9919\) trên tập nhà kính chưa từng xuất hiện lúc huấn luyện, vượt trội rõ rệt so với cả hai mô hình cơ sở B0 và B1. Toàn bộ 5 container được đóng gói bằng Docker Compose trong ngân sách bộ nhớ 1.792 MB, đáp ứng điều kiện triển khai thực tế trên một máy chủ VPS Linux phổ thông.

Bên cạnh những kết quả đạt được, nhóm cũng nhìn nhận thẳng thắn các hạn chế còn tồn tại trong phạm vi đồ án. Nguồn dữ liệu thu hoạch thực tế (\texttt{source = 'real'}) còn ít do chu kỳ sinh trưởng cây thật cần nhiều tuần và mô hình phần cứng mới có 2 vùng trồng, vì vậy mô hình học máy hiện vẫn dựa chủ yếu vào tập tổng hợp và tập thực nghiệm công khai WUR. Ngoài ra, việc sử dụng cảm biến môi trường phổ thông khiến hệ thống phải dựa trên các giả định xấp xỉ (\(T_{\text{lá}} = T_{\text{khí}} - 1\,^\circ\text{C}\) và quy đổi từ \texttt{lux} sang PPFD), đồng thời thang thời gian \texttt{demo} của hộp nhà kính mini cần được hiệu chỉnh lại bằng bộ tham số \texttt{real} khi áp dụng cho nhà kính thương mại. Ở khía cạnh hoàn thiện sản phẩm, còn hai cảnh báo tĩnh nhỏ từ \texttt{ruff} tại tệp \texttt{platform/api/routers/zones.py} và việc bổ sung nén thư mục \texttt{/models} cùng điều chỉnh \texttt{KEEP=7} trong kịch bản sao lưu \texttt{deploy/backup.sh} cần được dọn gọn trước khi nghiệm thu vận hành dài hạn.

Trong hướng phát triển tiếp theo, nhóm định hướng tiếp tục hiệu chuẩn bộ tham số \texttt{real} cho môi trường nhà kính sản xuất quy mô lớn và tối ưu cơ chế tiết kiệm năng lượng khi lấy mẫu ở tần suất thưa. Ở tầng nền tảng và học máy, hệ thống có thể mở rộng tính năng xuất nhật ký canh tác theo chuẩn VietGAP/GlobalGAP, tích hợp giải thích mức độ đóng góp đặc trưng bằng SHAP/\texttt{pred\_contrib} trực tiếp trên giao diện và phát triển cảnh báo bảo trì cảm biến sớm dựa trên xu hướng RSSI cũng như tỷ lệ mất gói khi tích lũy đủ dữ liệu thu hoạch thực địa từ nhiều trang trại.
